\section{Methodology}
\label{sec:method}

% We propose \textbf{TS-Memory}, a lightweight retrieval-free memory module for adapting frozen TSFMs to target domains. While online retrieval effectively mitigates distribution shifts, it introduces heavy inference-time search overhead. \textit{Our core insight is that retrieval not only provides point estimates but also induces a non-parametric conditional distribution over plausible futures.} Dispersion among retrieved neighbors inherently encodes context-dependent uncertainty, which is difficult to learn from a single realized trajectory. We treat this retrieval-induced distribution as \textbf{privileged offline supervision} and distill it into a parametric module, preserving retrieval-based adaptivity while enabling constant-time inference. As shown in Figure~\ref{fig:ts_plugmem_overview}, the framework operates in two stages.

We propose \textbf{TS-Memory}, a lightweight retrieval-free memory module for adapting frozen TSFMs to target domains. While online retrieval effectively mitigates distribution shifts, it introduces heavy inference-time search overhead. Our core insight is that retrieval not only provides point estimates but also induces a non-parametric conditional distribution over plausible futures. Dispersion among retrieved neighbors inherently encodes context-dependent uncertainty, difficult to learn from a single realized trajectory. We treat this retrieval-induced distribution as privileged offline supervision and distill it into a parametric module, preserving retrieval-based adaptivity while enabling constant-time inference. Since constructing such targets is leakage-sensitive and too costly for the online path, we adopt a two-stage design as shown in Figure~\ref{fig:ts_plugmem_overview}.


\begin{itemize}[leftmargin=*]
    \item \textbf{Stage I: Privileged Supervision Construction.} To prevent test-time leakage, we construct a distributional teacher exclusively from the training set. We retrieve $K$ nearest neighbors in the frozen embedding space, synthesize their future trajectories into empirical quantiles, and thereby mine a privileged supervision signal ($\mathcal{D}_{\text{teach}}$) that encodes local domain patterns and uncertainty structures without relying on a run-time index.
    
    \item \textbf{Stage II: Confidence-Gated Memory Distillation.} We train the memory module to directly predict teacher quantiles from raw context. To mitigate noisy retrieval, we design a \textit{confidence-gated} dual objective: the module aligns with ground truth via standard regression, and distills the retrieval-induced distribution only when the teacher provides a reliable high-confidence improvement signal. At inference, PlugMem runs in parallel with the frozen backbone for robust constant-time forecasting.
\end{itemize}

\subsection{Privileged Supervision Construction}
\label{sec:teacher}

TS-Memory treats retrieval as training-time privileged supervision used only offline to build an auxiliary teacher dataset:
\begin{equation}
\mathcal{D}_{\text{teach}}
=\left\{(\mathbf{X}_t,\mathbf{Y}_t,\widetilde{\mathbf{Q}}_t,\mathrm{Conf}_t)\right\},
\label{eq:tsm-dteach}
\end{equation}
% where $\widetilde{\mathbf{Q}}_t\in\mathbb{R}^{Q\times H\times C}$ are teacher quantile targets and $\mathrm{Conf}_t\in[0,1]$ measures the reliability of the retrieved evidence. 
Concretely, a leakage-safe kNN teacher has privileged access to a knowledge base of past contexts and their realized futures, which is unavailable to the student at test time. This access allows the teacher to form a retrieval-conditioned predictive distribution, summarized as quantile targets $\widetilde{\mathbf{Q}}_t\in\mathbb{R}^{Q\times H\times C}$ with reliability score $\mathrm{Conf}_t\in[0,1]$. The construction process involves three steps:

% The construction process involves three steps:% While each training window provides only one realized future $\mathbf{Y}_t$, retrieval provides a set of plausible futures from similar contexts; summarizing these futures into quantiles yields an instance-conditioned distributional target that complements point supervision.

\noindent\textbf{Knowledge Base Construction.} We first build a leakage-safe index utilizing the training split only: $\mathcal{K}=\{(\mathbf{X}^{(i)},\mathbf{Y}^{(i)})\}_{i=1}^{N_{\text{train}}}$,
ensuring that both the context window and forecast horizon are fully contained within the training segment. Then, we compute embeddings using a frozen encoder $f_{\text{enc}}$ (e.g., the encoder of the TSFM):
\begin{equation}
\mathbf{e}_t=f_{\text{enc}}(\mathbf{X}_t), \qquad\mathbf{e}_i=f_{\text{enc}}(\mathbf{X}^{(i)}).
\label{eq:tsm-emb}
\end{equation}
Candidates are retrieved via Euclidean distance $d_i=\|\mathbf{e}_t-\mathbf{e}_i\|_2$. We explicitly exclude the index-matched window to prevent trivial self-retrieval. Here, $\mathbf{e}_t$ is the embedding of the current input $\mathbf{X}_t$, and $\mathbf{e}_i$ is the embedding of the $i$-th candidate context. We use “retrieval-leakage-safe” to denote that all context–future pairs used by the teacher and retrieval baselines are constructed strictly from the training split; this does not assume that every TSFM pretraining corpus is fully auditable.

\noindent\textbf{Shift Alignment and Re-ranking.}
Distribution shift often manifests as level offsets, meaning that  neighbors close in embedding space may be mis-aligned in absolute scale. We align each candidate using a trailing mean shift over the last $m$ steps, computed channel-wise:
\begin{equation}
\boldsymbol{s}_i=\mathrm{mean}\!\left(\mathbf{X}_{t,L-m+1:L}\right)-\mathrm{mean}\!\left(\mathbf{X}^{(i)}_{L-m+1:L}\right)\in\mathbb{R}^{C}.
\label{eq:tsm-shift}
\end{equation}
This is applied to the context and future of the retrieved candidate:
\begin{equation}
\mathbf{X}^{(i)}_{\text{align}}=\mathbf{X}^{(i)}+\boldsymbol{s}_i,\qquad
\mathbf{Y}^{(i)}_{\text{align}}=\mathbf{Y}^{(i)}+\boldsymbol{s}_i.
\label{eq:tsm-align}
\end{equation}
Candidates are then re-ranked by the $\ell_1$ distance (i.e., $\text{score}_i=\|\mathbf{X}_t-\mathbf{X}^{(i)}_{\text{align}}\|_1$) between the query context $\mathbf{X}_t$ and the aligned candidate context $\mathbf{X}^{(i)}_{\text{align}}$, keeping the top $K$ neighbors.

\noindent\textbf{Teacher Aggregation and Confidence Scoring.}
Given the top $K$ neighbors with distances $\{d_k\}_{k=1}^{K}$, we compute softmax weights:
\begin{equation}
w_k=\frac{\exp\!\left(-\psi(d_k)/\tau_{\text{ret}}\right)}
{\sum_{j=1}^{K}\exp\!\left(-\psi(d_j)/\tau_{\text{ret}}\right)},
\qquad \sum_{k=1}^K w_k=1,
\label{eq:tsm-ret-weights}
\end{equation}
where $\psi(\cdot)$ is a monotone distance transform (identity by default) and $\tau_{\text{ret}}$ controls the weight sharpness. For each $u\in\mathcal{U}$, set $v_k=Y^{(k)}_{\text{align},u}$ and let $\pi$ sort $\{v_k\}$ in ascending order. Define cumulative weights $S_m=\sum_{r=1}^{m} w_{\pi(r)}$. Then, the weighted empirical quantile is:
\begin{equation}
\widetilde{Q}_{t,j,u}=v_{\pi(m^\star)},\qquad
m^\star=\min\{m:~S_m\ge q_j\}.
\label{eq:tsm-teacher-quantile}
\end{equation}
Here, $v_k$ is the aligned future value from the $k$-th neighbor, and $\pi(r)$ denotes the index of the $r$-th smallest value. We define retrieval confidence as the concentration of retrieval weights:
\begin{equation}
\mathrm{Conf}_t=\max_{1\le k\le K} w_k.
\label{eq:tsm-conf}
\end{equation}
The complete procedure is detailed in Algorithm~\ref{alg:teacher} (Appendix~\ref{appx:Offline_Teacher_Construction_Algorithm}).

\subsection{Confidence-Gated Memory Distillation}
\label{sec:plugmem}

We parameterize the memory module $g_{\phi}$ as a lightweight encoder--decoder Transformer. The input $\mathbf{X}_t$ is first normalized via Instance Normalization (Eq.~\ref{eq:tsm-inorm}), partitioned into patches of length $p$, and projected to $d$-dimensional tokens. A Transformer encoder produces a memory representation, which is attended to by $H$ learnable horizon queries in the decoder. A final quantile head maps these features to $Q$ quantile values, after which the normalization is inverted. Crucially, $g_{\phi}$ operates independently of the backbone internals, ensuring plug-and-play compatibility. 

Instance Normalization is defined as:
\begin{equation}
\widetilde{\mathbf{X}}_t = (\mathbf{X}_t - \mu(\mathbf{X}_t)) / (\sigma(\mathbf{X}_t) + \varepsilon).
\label{eq:tsm-inorm}
\end{equation}

To facilitate the memory module in learning transferable effective knowledge, we introduce three losses for joint optimization:

\noindent\textbf{Task Supervision $\mathcal{L}_{\text{task}}$.} Supervise PlugMem via standard quantile regression on ground-truth future values:
\begin{equation}
\mathcal{L}_{\text{task}} = \mathbb{E}_{t}\Big[
\operatorname{Pinball}_{\mathcal{Q}}\!\big(
\widehat{\mathbf{Q}}^{\text{mem}}_t,\mathbf{Y}_t
\big) \Big],
\label{eq:task-pinball-abstract}
\end{equation}
where $\operatorname{Pinball}_{\mathcal{Q}}(\cdot,\cdot)$ denotes the standard pinball loss for quantile regression, aggregated over the forecast horizon and variables.

\noindent\textbf{Confidence-Gated Distillation $\mathcal{L}_{\text{align}}$.} Distill the teacher's adaptive improvement into the memory module via confidence-gated distillation (two steps: error evaluation and incremental alignment), activated only when retrieval provides a reliable signal:

First, we evaluate the teacher's absolute prediction quality via median absolute error (median index $j^\star=\arg\min_j |q_j-0.5|$) to screen valid distillation windows, defining median absolute error for the retrieval teacher ($\mathrm{err}_t^T$) and frozen backbone ($\mathrm{err}_t^\mathrm{base}$) as:
\begin{equation}
\mathrm{err}_t^\star=\frac{1}{|\mathcal{U}|}\sum_{u\in\mathcal{U}}\bigl|Y_{t,u}-Q_{t,j^\star,u}^\star\bigr|,\;\star\in\{T,\mathrm{base}\},
\label{equ:err}
\end{equation}
where $\mathrm{err}_t^T$ and $\mathrm{err}_t^\mathrm{base}$ quantify the absolute deviation of the teacher and backbone predictions from the ground truth $Y_{t,u}$, respectively.

We then gate distillation with a margin $\epsilon_{\text{gate}}$ to retain only windows where the teacher outperforms the backbone:
\begin{equation}
\chi_t=\mathbb{I}\!\left(\mathrm{err}^{T}_t+\epsilon_{\text{gate}}<\mathrm{err}^{\text{base}}_t\right),\quad
\omega_t=\chi_t\cdot \mathrm{Conf}_t^{\gamma},
\label{eq:tsm-gate}
\end{equation}
where $\chi_t$ is the binary gating indicator (1 for valid windows), $\omega_t$ the confidence-weighted gating coefficient, $\mathrm{Conf}_t$ the retrieval confidence score, and $\gamma$ the confidence scaling hyperparameter. 



\begin{table*}[!htbp]
\begin{center}
\caption{Long-term forecasting results of TS-Memory across frozen TSFM backbones. Results are averaged over forecasting horizons $H \in $\{96, 192, 336, 720\}. Lower values indicate better performance. Best results are highlighted in \textbf{bold}, and second
best results are underlined. Values in parentheses denote relative
improvement (\%). Full results are in Table~\ref{tab:ts_plugmem_full} of Appendix~\ref{appx:complete_results}.}

\begin{small}
\setlength\tabcolsep{3pt}
\renewcommand{\arraystretch}{0.85} 
\scalebox{0.9}{
\begin{tabular}{c|cc|cc|cc|cc|cc|cc|cc|cc}
\toprule
\multirow{1}{*}{TSFM} &
\multicolumn{4}{c|}{ChronosBolt} &
\multicolumn{4}{c|}{Chronos2} &
\multicolumn{4}{c|}{Sundial} &
\multicolumn{4}{c}{TimesFM} \\
\cmidrule{1-17}
\multirow{1}{*}{Dataset}
 &
\multicolumn{2}{c|}{Origin} & \multicolumn{2}{c|}{TS-Memory} &
\multicolumn{2}{c|}{Origin} & \multicolumn{2}{c|}{TS-Memory} &
\multicolumn{2}{c|}{Origin} & \multicolumn{2}{c|}{TS-Memory} &
\multicolumn{2}{c|}{Origin} & \multicolumn{2}{c}{TS-Memory} \\
\midrule
Metric &
MSE & MAE & MSE & MAE &
MSE & MAE & MSE & MAE &
MSE & MAE & MSE & MAE &
MSE & MAE & MSE & MAE \\
\midrule

ETTh1 & \underline{0.448} & \underline{0.419} & \textbf{0.421}\pct{-6.0\%} & \textbf{0.414}\pct{-1.2\%} & \underline{0.442} & \underline{0.412} & \textbf{0.420}\pct{-5.0\%} & \textbf{0.407}\pct{-1.2\%} & \underline{0.400} & \underline{0.412} & \textbf{0.396}\pct{-0.9\%} & \textbf{0.410}\pct{-0.5\%} & \underline{0.479} & \underline{0.442} & \textbf{0.448}\pct{-6.5\%} & \textbf{0.431}\pct{-2.5\%} \\
\midrule
ETTh2 & \underline{0.367} & \underline{0.380} & \textbf{0.354}\pct{-3.4\%} & \textbf{0.375}\pct{-1.1\%} & \underline{0.376} & \underline{0.383} & \textbf{0.350}\pct{-6.9\%} & \textbf{0.378}\pct{-1.5\%} & \underline{0.344} & \underline{0.380} & \textbf{0.340}\pct{-1.1\%} & \textbf{0.378}\pct{-0.6\%} & \underline{0.402} & \underline{0.409} & \textbf{0.364}\pct{-9.4\%} & \textbf{0.398}\pct{-2.9\%} \\
\midrule
ETTm1 & \underline{0.421} & \underline{0.383} & \textbf{0.381}\pct{-9.4\%} & \textbf{0.371}\pct{-3.1\%} & \underline{0.433} & \underline{0.381} & \textbf{0.393}\pct{-9.2\%} & \textbf{0.372}\pct{-2.3\%} & \underline{0.369} & \underline{0.369} & \textbf{0.356}\pct{-3.6\%} & \textbf{0.364}\pct{-1.5\%} & \underline{0.429} & \underline{0.416} & \textbf{0.382}\pct{-11.0\%} & \textbf{0.396}\pct{-4.8\%} \\
\midrule
ETTm2 & \underline{0.291} & \underline{0.317} & \textbf{0.267}\pct{-8.0\%} & \textbf{0.311}\pct{-2.1\%} & \underline{0.295} & \underline{0.315} & \textbf{0.270}\pct{-8.4\%} & \textbf{0.309}\pct{-2.0\%} & \underline{0.276} & \underline{0.317} & \textbf{0.266}\pct{-3.6\%} & \textbf{0.312}\pct{-1.6\%} & \underline{0.332} & \underline{0.341} & \textbf{0.279}\pct{-16.0\%} & \textbf{0.324}\pct{-5.0\%} \\
\midrule
Electricity & \underline{0.159} & \underline{0.244} & \textbf{0.155}\pct{-2.5\%} & \textbf{0.241}\pct{-1.0\%} & \underline{0.163} & \underline{0.244} & \textbf{0.159}\pct{-2.1\%} & \textbf{0.241}\pct{-1.1\%} & \underline{0.148} & \underline{0.242} & \textbf{0.144}\pct{-2.7\%} & \textbf{0.239}\pct{-1.4\%} & \underline{0.154} & \underline{0.244} & \textbf{0.149}\pct{-2.9\%} & \textbf{0.241}\pct{-1.1\%} \\
\midrule
Exchange & \underline{0.371} & \underline{0.412} & \textbf{0.365}\pct{-1.8\%} & \textbf{0.406}\pct{-1.5\%} & \underline{0.399} & \underline{0.421} & \textbf{0.366}\pct{-8.3\%} & \textbf{0.408}\pct{-3.1\%} & \underline{0.553} & \underline{0.494} & \textbf{0.471}\pct{-14.8\%} & \textbf{0.464}\pct{-6.0\%} & \underline{0.433} & \underline{0.446} & \textbf{0.425}\pct{-1.9\%} & \textbf{0.433}\pct{-2.7\%} \\
\midrule
Traffic & \underline{0.435} & \underline{0.263} & \textbf{0.425}\pct{-2.3\%} & \textbf{0.258}\pct{-1.6\%} & \underline{0.394} & \underline{0.237} & \textbf{0.389}\pct{-1.4\%} & \textbf{0.235}\pct{-0.6\%} & \underline{0.461} & \underline{0.286} & \textbf{0.436}\pct{-5.3\%} & \textbf{0.278}\pct{-3.0\%} & \underline{0.370} & \underline{0.244} & \textbf{0.367}\pct{-0.6\%} & \textbf{0.241}\pct{-1.1\%} \\
\midrule
Weather & \underline{0.263} & \underline{0.276} & \textbf{0.239}\pct{-9.1\%} & \textbf{0.267}\pct{-3.5\%} & \underline{0.274} & \underline{0.267} & \textbf{0.241}\pct{-12.0\%} & \textbf{0.260}\pct{-2.6\%} & \underline{0.244} & \underline{0.269} & \textbf{0.238}\pct{-2.5\%} & \textbf{0.266}\pct{-1.2\%} & \underline{0.222} & \underline{0.237} & \textbf{0.205}\pct{-7.6\%} & \textbf{0.230}\pct{-2.9\%} \\

\bottomrule
\end{tabular}}
\end{small}
\label{tab:ts_plugmem_compact}
\end{center}

\end{table*}



Second, we perform incremental alignment to distill the teacher's adaptive improvement (not global bias) into the memory module, aligning quantile predictions via the robust Huber loss $\ell_{\kappa}(\cdot,\cdot)$:
\begin{equation}
\mathcal{D}_{Q}(t) = \frac{1}{Q|\mathcal{U}|}\sum_{j=1}^{Q}\sum_{u\in\mathcal{U}}\ell_{\kappa}\!\left(\widehat{Q}^{\text{mem}}_{t,j,u},\widetilde{Q}_{t,j,u}\right).
\label{eq:tsm-dq}
\end{equation}
To mitigate sensitivity to global central bias of the frozen backbone, we further align the incremental median correction between the teacher and memory module (relative to the backbone):
\begin{equation}
\Delta^{\text{mem}}_{t,u}=\widehat{Q}^{\text{mem}}_{t,j^\star,u}-\widehat{Q}^{\text{base}}_{t,j^\star,u}, \quad
\qquad
\Delta^{T}_{t,u}=\widetilde{Q}_{t,j^\star,u}-\widehat{Q}^{\text{base}}_{t,j^\star,u},
\label{eq:tsm-delta}
\end{equation}
where $\Delta^{\text{mem}}_{t,u}$ and $\Delta^{T}_{t,u}$ denote the memory module and teacher's incremental median correction relative to the frozen backbone, respectively. The alignment loss for these corrections is:
\begin{equation}
\mathcal{D}_{\Delta}(t) = \frac{1}{|\mathcal{U}|}\sum_{u\in\mathcal{U}}
\ell_{\kappa}\!\left(\Delta^{\text{mem}}_{t,u},~\Delta^{T}_{t,u}\right).
\label{eq:tsm-ddelta}
\end{equation}
The final confidence-gated alignment loss is:
\begin{equation}
\mathcal{L}_{\text{align}}=\mathbb{E}_{t}\!\left[\omega_t\cdot\bigl(\mathcal{D}_{Q}(t)+\eta\,\mathcal{D}_{\Delta}(t)\bigr)\right].
\label{eq:tsm-loss-align}
\end{equation}




\noindent\textbf{Stability Regularization $\mathcal{L}_{\text{reg}}$.} Constrain the memory module to conservative behavior for uncertain retrieval (small $\omega_t$, weak $\mathcal{L}_{\text{align}}$ supervision), avoiding over-adjustment of the frozen backbone. We first anchor its median to the backbone with weight $(1-\omega_t)$:
\begin{equation}
\mathcal{L}_{\text{anchor}} = \mathbb{E}_{t}\!\left[(1-\omega_t)\cdot\frac{1}{|\mathcal{U}|}\sum_{u\in\mathcal{U}}\ell_{\kappa}\!\left(\widehat{Q}^{\text{mem}}_{t,j^\star,u},\widehat{Q}^{\text{base}}_{t,j^\star,u}\right)\right].
\label{eq:tsm-loss-anchor}
\end{equation}
We additionally impose a universal regularization to prevent quantile crossing in memory predictions, ensuring monotonic outputs:
\begin{equation}
\mathcal{L}_{\text{cross}} = \mathbb{E}_{t}\!\left[\frac{1}{(Q-1)|\mathcal{U}|}\sum_{u\in\mathcal{U}}\sum_{j=1}^{Q-1}\max\!\left(0,\widehat{Q}^{\text{mem}}_{t,j,u}-\widehat{Q}^{\text{mem}}_{t,j+1,u}\right)\right].
\label{eq:tsm-loss-cross}
\end{equation}
The final stability regularization loss is: $\mathcal{L}_{\text{reg}}=\mathcal{L}_{\text{anchor}}+\lambda_{\text{cross}}\,\mathcal{L}_{\text{cross}}$.

\noindent\textbf{Training objective.}
We train the memory module on $\mathcal{D}_{\text{teach}}$ using a composite loss that balances ground-truth regression with distributional distillation:
\begin{equation}
\mathcal{L} = \mathcal{L}_{\text{task}} + \lambda_{\text{align}}\,\mathcal{L}_{\text{align}} + \lambda_{\text{reg}}\,\mathcal{L}_{\text{reg}}.
\label{eq:tsm-loss-total}
\end{equation}

\subsection{Inference via Adaptive Fusion}
\label{sec:inference}

At test time, TS-Memory performs strictly retrieval-free inference:
\begin{equation}
\widehat{\mathbf{Q}}^{\text{base}}_t=f_{\theta}(\mathbf{X}_t),
\qquad
\widehat{\mathbf{Q}}^{\text{mem}}_t=g_{\phi}(\mathbf{X}_t).
\label{eq:tsm-infer}
\end{equation}
We fuse the two quantile forecasts by quantile-wise interpolation:
\begin{equation}
\widehat{\mathbf{Q}}^{\text{final}}_t
=(1-\alpha)\widehat{\mathbf{Q}}^{\text{base}}_t+\alpha\widehat{\mathbf{Q}}^{\text{mem}}_t,
\qquad
\alpha\in[0,1].
\label{eq:tsm-fuse}
\end{equation}
We tune $\alpha$ on a validation split.
Inference requires only two forward passes: one through the frozen backbone and one through the memory module.
No retrieval index needs to be maintained and no nearest-neighbor search is performed at serving time.

\noindent\textbf{Inference-time complexity.} TS-Memory adds constant overhead: one lightweight extra forward pass and quantile-wise fusion. Runtime and memory do not depend on the knowledge base since retrieval is offline. In contrast, online retrieval requires query-time search and aggregation that scale with the datastore.
